% Hints for VI/03. Not included in the default student build.

\chapter*{Hints for VI/03}
\addcontentsline{toc}{chapter}{Hints for VI/03}

\paragraph{Exercise \ref{ex:vi03-kernel}.}
Unpack what it means for a restricted polynomial function to be the zero function.

\paragraph{Exercise \ref{ex:vi03-same-function}.}
Apply the previous exercise to \(f-g\).

\paragraph{Exercise \ref{ex:vi03-point}.}
Use
\[
I(\{a\})=(x_1-a_1,\dots,x_n-a_n)
\]
and evaluation at \(a\).

\paragraph{Exercise \ref{ex:vi03-line}.}
Solve the equation for \(y\) because \(b\neq0\).

\paragraph{Exercise \ref{ex:vi03-parabola}.}
Send \(\bar x\mapsto t\), \(\bar y\mapsto t^2\), and send \(t\mapsto\bar x\) in the other
direction.

\paragraph{Exercise \ref{ex:vi03-hyperbola}.}
Inside the quotient, \(\bar x\bar y=1\).

\paragraph{Exercise \ref{ex:vi03-reduced}.}
Use that \(I(X)\) is radical.

\paragraph{Exercise \ref{ex:vi03-generators}.}
Every residue class is represented by a polynomial in the ambient variables.

\paragraph{Exercise \ref{ex:vi03-evaluation}.}
Constants show surjectivity.  A quotient by the kernel is \(k\).

\paragraph{Exercise \ref{ex:vi03-closed-subset}.}
Apply the third isomorphism theorem to
\[
I(X)\subseteq I(Y).
\]

\paragraph{Exercise \ref{ex:vi03-two-points}.}
Use
\[
I(\{a,b\})=((x-a)(x-b))
\]
and the Chinese remainder theorem.

\paragraph{Exercise \ref{ex:vi03-finite-points}.}
The ideals \((x-a_i)\) are pairwise comaximal.

\paragraph{Exercise \ref{ex:vi03-zero-divisors}.}
The classes of \(x\) and \(y\) multiply to zero.

\paragraph{Exercise \ref{ex:vi03-nilpotent-warning}.}
Look at the class of \(x\) in \(B\).

\paragraph{Exercise \ref{ex:vi03-finite-field}.}
First show that \(x^q-x\) vanishes everywhere.  Divide an arbitrary vanishing polynomial by
\(x^q-x\).

\paragraph{Exercise \ref{ex:vi03-challenge-circle}.}
Compute
\[
(\bar x+i\bar y)(\bar x-i\bar y).
\]
